\documentclass[12pt, a4paper]{article}

\usepackage{graphicx}
\usepackage{amsmath, amssymb}
\usepackage[T2A]{fontenc}
\usepackage[english, russian]{babel}
\usepackage[utf8]{inputenc}
\usepackage[margin=2cm]{geometry}
\usepackage[most]{tcolorbox}
\usepackage{caption}
\usepackage{enumitem}

\tcbuselibrary{breakable}
\tcbset{
  width=0.9\textwidth,
  halign=justify,
  center,
  breakable,
  colback=white
}

\title{Конспект по математическому анализу}
\author{VG6, Лектор Теляковский Д.С.}
\date{2 семестр 2025/2026}

\newcommand{\N}{\mathbb{N}}
\newcommand{\Q}{\mathbb{Q}}
\newcommand{\Z}{\mathbb{Z}}
\newcommand{\R}{\mathbb{R}}
\newcommand{\eps}{\varepsilon}

\newcommand{\sign}{\text{sign\ }}

\begin{document}
\maketitle
\tableofcontents
\newpage

\section{Исследование функций}
\subsection{Условия постоянства и монотонности.}

\textbf{Теорема.}
\begin{tcolorbox}
	\begin{enumerate}
		\item $f`(x) > 0 \Rightarrow f(x) \text{строго возрастает} \Rightarrow f`(x) \geq 0$ 
		\item $f`(x) \geq 0 \Rightarrow f(x) \text{возрастает} \Rightarrow f`(x) \geq 0$ 
		\item $f`(x) \equiv 0 \Rightarrow f(x) \equiv Const \Rightarrow f`(x) \equiv 0$ 
		\item $f`(x) \leq 0 \Rightarrow f(x) \text{убывает} \Rightarrow f`(x) \leq 0$ 
		\item $f`(x) < 0 \Rightarrow f(x) \text{строго убывает} \Rightarrow f`(x) \leq 0$ 
	\end{enumerate}
\end{tcolorbox}
\begin{tcolorbox}[title=Доказательство]
	\begin{enumerate}
		\item Левый столбйц. Формула Лагранжа конечных приращений. 
			\[ \forall x_1, x_2 \in (a, b), x_1 < x_2\; \exists \xi \in (x_1, x_2) \]
		\[ f(x_2) - f(x_1) = f`(\xi) (x_2 - x_1), \;\;\;\;\; (x_2 - x_1) > 0 \]
	\item Правый столбец. 
		\[ x \in (a, b),\; \forall h: \in (a, b)\;\frac{f(x+h) - f(x)}{h} \xrightarrow[]{h\to 0} f`(x)\]
	\end{enumerate}
\end{tcolorbox}

\textbf{Замечание.} $y = x^3$ строго возрастает. 
$y` = 3x^2$, при $x = 0$

\subsection{Условие внутреннего экстремума.}
\textbf{Теорема 1. Теорема Ферма.} $f(x)$ на $(a, b)$, $f(x)$ экстремум в точке $\xi \in (a, b) \Rightarrow f`(\xi) = 0$ или не существует. 
% нарисовать график y = x^3. Показать производную в 0. 
\textbf{Теорема 2.} $f(x) \in C(a, b), x_0 \in (a, b), f(x) \text{ дифференцируема на }(a, b)\setminus \{x_0\}$
\begin{enumerate}
	\item $f`(x) < 0$ на $(a, x_0), (x_0, b) \Rightarrow$ в $x_0$ нет экстремума. 
	\item $f`(x) < 0$ на $(a, x_0), f`(x) > 0$ на $(x_0, b) \Rightarrow$ в $x_0$ у $f(x)$ - минимум.
	\item $f`(x) > 0$ на $(a, x_0), (x_0, b) \Rightarrow$ в $x_0$ нет экстремума. 
	\item $f`(x) > 0$ на $(a, x_0), f`(x) < 0$ на $(x_0, b) \Rightarrow$ в $x_0$ у $f(x)$ - максимум.
\end{enumerate}
\begin{tcolorbox}[title=Доказательство]
	\begin{enumerate}
		\item $f`(x) < 0$ на $(a, x_0), (x_0, b)$ по утверждению $\forall x` \in (a, x_0)$ и $x`` \in (x_0, b)$
			\[ f(x`) > f(x) > f(x``) \Rightarrow \text{ в } x_0\text{ нет экстремума} \]
		\item $f`(x) < 0$ на $(a, x_0)$ и $f(x) > 0$ на $(x_0, b)$ по утверждению $\forall x` \in (a, x_0)$ и $x`` \in (x_0, b)$
			\[ f(x`) > f(x_0)\text{ и } f(x_0) < f(x``) \]
	\end{enumerate}
\end{tcolorbox}


\textbf{Утверждение.} $f(x) \in C([a, b])$, дифференцируема на $(a, b)$
\begin{enumerate}
	\item $f`(x) > 0$ на $(a, b) \Rightarrow \min_{[a, b]} f(x) = f(a), \max_{[a, b]} f(x) = f(b)$
	\item $f`(x) < 0$ на $(a, b) \Rightarrow \max_{[a, b]} f(x) = f(a), \min_{[a, b]} f(x) = f(b)$
\end{enumerate}
\begin{tcolorbox}[title=Доказательство.]
	\begin{enumerate}
		\item $\forall x_0 \in (a, b)$ Для $f(x)$ на $[a, x_0]$ и $[x_0, b]$ то по Т. Лагранжа о конечных приращениях $\exists \xi` \in (x, x_0), \xi`` \in (x_0, b)$
			\[ f(x_0) - f(a) =f`(\xi`)(x_0 - a) > 0 \]
			\[ f(b) - f(x_0) =f`(\xi``)(b - x_0) > 0\]
		\begin{enumerate}
			\item $f`(x) > 0$ на $(a, b) \Rightarrow f(x) < f(x_0) < f(b)$
		\end{enumerate}
	\end{enumerate}
\end{tcolorbox}

\textbf{Теорема 3. Достаточное условие в теореме высших производных.} 
$f(x)$ в $U(x_0)$ и имеет в $x_0$ производную порядка $n \geq 2$, и выполняется:
\[ f`(x_0) = f``(x_0) = ...  = f^{(n - 1)}(x_0), f^{(n)}(x_0) \not= 0\]
Тогда, если 
\begin{enumerate}
	\item n - нечётное, то в $x_0$ экстремума нет. 
	\item n - чётное, то экстремум есть. 
		\begin{enumerate}
			\item $f^{(n)}(x_0) > 0 \Rightarrow$ в $x_0$ локальный минимум
			\item $f^{(n)}(x_0) < 0 \Rightarrow$ в $x_0$ локальный максимум
		\end{enumerate}
\end{enumerate}
\begin{tcolorbox}[title=Доказательство Формулы Тейлора и остатком по Пиано.]
	Используем формулу Тейлора с остаточным членом в форме Пеано:
	\[ f(x) = f(x_0) + \frac{f'(x_0)}{1!}(x-x_0) + \dots\] 
	\[... + \frac{f^{(n-1)}(x_0)}{(n-1)!}(x-x_0)^{n-1} + \frac{f^{(n)}(x_0)}{n!}(x-x_0)^n + \bar{o}((x-x_0)^n) \]
	
	Так как $f'(x_0) = \dots = f^{(n-1)}(x_0) = 0$, выражение упрощается:
	\[ f(x) = f(x_0) + \frac{f^{(n)}(x_0)}{n!}(x-x_0)^n + \bar{o}((x-x_0)^n) \]
	
	Распишем $\bar{o}$ как бесконечно малую функцию $\alpha(x)$:
	\[ f(x) = f(x_0) + \frac{f^{(n)}(x_0)}{n!}(x-x_0)^n + \alpha(x)(x-x_0)^n, \quad \text{где } \alpha(x) \xrightarrow{x \to x_0} 0 \]
	
	Вынесем $(x-x_0)^n$ за скобку:
	\[ f(x) = f(x_0) + \left( \frac{f^{(n)}(x_0)}{n!} + \alpha(x) \right) (x-x_0)^n \]
	
	Так как $\alpha(x) \to 0$, то для члена в скобках справедливо:
	\[ \exists V(x_0): \forall x \in V(x_0) \quad \sign \left( \frac{f^{(n)}(x_0)}{n!} + \alpha(x) \right) = \sign f^{(n)}(x_0) \]
	
	Таким образом, для $\forall x \in V(x_0)$:
	\[ f(x) = f(x_0) + \underbrace{\left( \frac{f^{(n)}(x_0)}{n!} + \alpha(x) \right)}_{\text{сохр. знак}} (x-x_0)^n \]
	
\end{tcolorbox}

\subsection{Выпуклые функции.}
\textbf{Определение 1.} $f(x)$ на $<a, b>$ \textbf{выпукла вниз}, если $\forall x_1, x_2 \in <a, b>$ график функции $y = f(x)$ на $(x_1, x_2)$ раположена нижу отрезка с концами в $(x_1, f(x_1)), (x_2, f(x_2))$. (строгая выпуклость)

% Нарисованить график выпуклой вниз функции с отрезком 

Вопрос: Пусть функция выпукла вниз а интервале. Доказатель, что она непрерывна на всём интервале. 

\textbf{Нестрогая выпуклость} - не выше. 

% тожу нужен рисунок

\[ \forall x \in (x_1, x_2):\; y_{\text{сек}} = \boxed{\frac{x_2 - x}{x_2 - x_1}f(x_1) + \frac{x - x_1}{x_2 - x_1}f(x_2) \geq f(x) }\]
\[
\left( \frac{x_2 - x}{x_2 - x_1} + \frac{x - x_1}{x_2 - x_1} \right) f(x) \leq \frac{x_2 - x}{x_2 - x_1} f(x_1) + \frac{x - x_1}{x_2 - x_1} f(x_2)
\]

\[
\frac{x_2 - x}{x_2 - x_1} \left( f(x) - f(x_1) \right) \leq \frac{x - x_1}{x_2 - x_1} \left( f(x_2) - f(x) \right)
\]

\[
\boxed{ \frac{f(x) - f(x_1)}{x - x_1} \leq \frac{f(x_2) - f(x)}{x_2 - x} }
\]
1. Что будет если функция дифференцируема
3. Доказать что в каждой точке существует односторонний придел, даже если она не дифференцируема. 
\end{document}
